
\chapter{引言}
粒子物理学是研究宇宙中最基本的构成元素及其性质和相互作用的基础物理学。目前描述基本粒子最成功的模型是
标准模型。该模型假定粒子的本质是场，如电子场，夸克场，光子场等等，然后再对场进行量子化。这种把物质场
进行量子化的理论通称为量子场论。

人们对物质世界量子性质的探索始于二十世纪初。到1925年，经过约二十年的发展，量子力学已基本建立。
1926年，M.~Born、W.~Heisenberg和P.~Jordan等人开始尝试对经典电磁场进行量子化。
此后，人们利用正则量子化的方法处理各种经典场（历史上也称为二次量子化）。但量子场论给出的理论
预言却到处都是发散（无穷大）。这些发散可以归为两类：紫外发散和红外发散。本质上来说，所有的发散都源于无穷多
自由度体系的假定。发散问题一直困扰着理论物理学家们，其实不止在量子场论中，在经典电动力学中如电子自能都存在
发散。后来J.~Schwinger，R.~Feynman和Tomonaga发现紫外发散可以吸收到理论中的参数里面。随后1949年，F.~Dyson
证明了在量子电动力学（QED）中，通过参数的重新定义可以消除紫外发散，所以此时QED的理论已经基本完备了。

%除了电磁相互作用外，人们很早的时候也发现了其它相互作用。
在核物理的现象中，人们发现放射性元素可以衰变出$\beta$射线（电子束）。对$\beta$衰变的研究，最成功的
早期模型是费米的四费米子理论。虽然四费米子理论能够成功解释当时的弱相互作用现象，但该理论不可重整化，
因此，人们继续寻找更深层次的物理规律。二十世纪五十年代，李政道与杨振宁提出弱相互作用中宇称不守恒的假设，
随后吴健雄通过实验证实了这一假设。由此，人们确认弱相互作用具有$V-A$结构。
1958年，R.~Feynman和Gell-Mann仿照QED提出了弱相互作用的理论描述。随后，S.~Glashow提出了弱电统一模型，
但该理论当时仍存在一些问题，例如不知道如何正确地给W、Z玻色子赋予质量。

关于原子核之间的强相互作用的研究始于H.~Yukawa的理论，Yukawa认为核子之间是通过交换$\pi$介子产生相互作用。
随着实验物理学的发展，$\pi$介子在宇宙线实验中被发现。后来，人们又陆续发现了许多强子。为了解释这些新发现
的强子，Gell-Mann等人提出了夸克模型\cite{QuarkModel}。1965年，O.~Greenberg，M.~Han和Y.~Nambu等人提出
夸克带有新的内禀自由度，也就是色荷。
介子和重子都是色单态粒子。由于实验上从来没有找到孤立的自由夸克，所以人们提出了色禁闭猜想，
即实验上探测到的粒子必须是色单态粒子。到了七十年代初，H.~Fritzsch和Gell-Mann在Yang-Mills场的基础上提出了
量子色动力学（QCD）的拉氏量。同一年，D.~Politzer，D.~Gross和F.~Wilczek提出了QCD的渐进自由性质
\cite{Politzer,GrossWilczek}，这可能是
导致夸克色禁闭的原因。后来在SLAC的实验证明了QCD的确有渐进自由的性质，
自此强相互作用的基本规律也确立了起来。

随着对强，弱和电磁相互作用的深入理解，物理学家认为弱相互作用也应该纳入规范场论当中。早在上个世纪六十年代，
Higgs，Kibble，Guralnik，Hagen，Brout和Englert等提出了非阿贝尔场下的自发对称破缺机制，也就是著名的Higgs
机制\cite{Higgs:1,Higgs:2,Higgs:3}。
有了这个机制，S.~Weinberg和A.~Salam在1967年提出Weinberg-Salam模型统一描述弱电相互作用
\cite{GWS:1,GWS:2,GWS:3}。
该模型的对称群是$SU(2)_L\otimes U(1)_Y$，并且预言了存在三个有质量的规范玻色子，即今天所说的$W^\pm$和Z
玻色子。再把QCD加进去，就构成了今天的标准模型。可以说，标准模型是现代物理学最成功的理论模型之一。
检验标准模型和发现新的基本粒子一直是粒子物理学家的梦想。1995年以前，除了Higgs粒子，其它所有标准模型粒子
都陆续在高能物理对撞机上发现。直到2012年，标准模型中最重要的粒子-Higgs玻色子才在位于欧洲的大型强子对撞机
（LHC）上发现\cite{ATLAS-2012,CMS-2012}。
此后，标准模型的预言与实验结果符合得很好。进一步检验Higgs粒子的性质仍然十分必要。

标准模型虽然取得了巨大成功，但它并非物理学的终极理论。一个更完整的理论还应该能够解释
引力，以及暗物质，暗能量和中微子质量等等。人们在标准模型的基础上提出了种种新模型，期望发现新物理。
目前人们提出了各种新物理模型，如超对称模型，额外维模型，Little Higgs模型，双Higgs二重态模型（2HDM）等等。
其中，2HDM是对标准模型Higgs部分的简单扩展，为研究额外Higgs粒子的存在提供了框架。本文主要研究
大型强子对撞机上WW+jet产生过程，以及2HDM中VBF道产生Higgs对过程的高阶修正。其中，WW+jet产生既可以用来研究
W玻色子的规范耦合，也可以作为新物理过程的背景。2HDM模型下VBF道产生Higgs对可以用来研究Higgs自耦合，也可以用来
探索新的Higgs粒子。

本论文的组织结构如下：第二章介绍标准模型、双Higgs二重态模型及电弱修正的相关理论；第三章讨论LHC上
WW+jet产生过程的次领头阶QCD和EW修正；第四章研究2HDM中VBF道Higgs对产生过程的NNLO QCD修正；
第五章总结全文。
